\scanpage{140}
% archive foliation: 139
2) Soit $S'\mathbin{\text{---}}S$ fid. plt et qu. cpt. Alors

i)
\[
\mathrm{Pic}'(X/S)\mathbin{\text{---}}\mathrm{Pic}'(X'/S')
\]
est injectif. ii) Si $X/\Sigma$ a \deleted{\unsure{une}} une section \deleted{sur} localement sur $S$, alors
\diagram{\begin{tikzcd}
\mathrm{Pic}'(X/S)\arrow[r]&\mathrm{Pic}'(X'/S')\arrow[r,shift left]\arrow[r,shift right]&\mathrm{Pic}'(X''/S'')
\end{tikzcd}}
est exact.
\margin{généraliser à des familles \unsure{recouvrantes} du type $\unsure{U^*}$ dès l'abord \unsure{---}}

On peut dém. (i) partant d'un $\uline{L}$ inv. sur $X$. \inserted{En vertu de 1)} L'hyp. signifie que $\uline{L}'$ provient de $\Sigma'$. Or $\Sigma'$ est fid. plt qu. cpt sur $\Sigma$. On sait que alors $\uline{L}$ provient de $\Sigma$.

Pour (ii), \inserted{grâce à (i)} la question est locale sur $S$, on peut donc supposer que $X$ a une section sur $\Sigma$. Alors on utilise le théorème de descente connu et le raisonnement connu.

\uline{Corollaire} \deleted{\unsure{1}} Dans les conditions précédentes, $\mathrm{Pic}'(X/S)\longrightarrow\mathrm{Pic}^{\unsure{\sharp|*}}(X/S)$ est injectif, et est bijectif si $X/\Sigma$ \deleted{\unsure{localement}} a une section loc. sur $S$.

Si $X/\Sigma$ a \unsure{=} une section, on a un isom
\[
\mathrm{Pic}(X/S)\simeq\mathrm{Pic}(X)/\mathrm{Pic}(\Sigma)
\]

